\documentclass[11pt]{article}

\usepackage[T1]{fontenc}
\usepackage{amsmath,amssymb,amsthm}
\usepackage{mathptmx}
\usepackage[margin=1in]{geometry}
\usepackage[authoryear,round]{natbib}
\usepackage{booktabs}
\usepackage{tikz}
\usepackage{pgfplots}
\pgfplotsset{compat=1.18}
\usepgfplotslibrary{groupplots}
\usepackage{caption}
\usepackage[hidelinks]{hyperref}
\hypersetup{pdftitle={Pricing Uncollateralized Microcredit on a Public Blockchain: Short Loans, Re-lending and a Locked First-Loss Reserve}, pdfauthor={Claude Code}}
\usepackage{xcolor}

\theoremstyle{plain}
\newtheorem{proposition}{Proposition}
\newtheorem{theorem}{Theorem}
\newtheorem{corollary}{Corollary}
\theoremstyle{definition}
\newtheorem{remark}{Remark}

\makeatletter
\newcommand{\inputrows}[1]{\@@input #1 }
\makeatother
\newcommand{\E}{\mathbb{E}}
\newcommand{\EL}{\mathrm{EL}}
\newcommand{\UL}{\mathrm{UL}}

% Generated by scripts/locked_reserve.py; do not edit.
\newcommand{\RElThree}{30}
\newcommand{\AprThree}{10.33}
\newcommand{\LongRunLockedThree}{3.95}
\newcommand{\LongRunAtElThree}{6.18}
\newcommand{\RRecThree}{55}
\newcommand{\LongRunFortyFiveThree}{4.83}
\newcommand{\RElFive}{42}
\newcommand{\AprFive}{12.33}
\newcommand{\LongRunLockedFive}{3.67}
\newcommand{\LongRunAtElFive}{6.10}
\newcommand{\RRecFive}{65}
\newcommand{\LongRunFortyFiveFive}{5.76}
\newcommand{\RElTen}{58}
\newcommand{\AprTen}{18.33}
\newcommand{\LongRunLockedTen}{3.90}
\newcommand{\LongRunAtElTen}{6.61}
\newcommand{\RRecTen}{75}
\newcommand{\LongRunFortyFiveTen}{6.61}
\newcommand{\Effr}{4.33}

% Generated by scripts/event_timed.py; do not edit.
\newcommand{\EventQHundred}{21.00}
\newcommand{\CombinedQHundred}{21.85}
\newcommand{\CombinedErrHundred}{+4.0}
\newcommand{\EventMeanLossHundred}{5.12}
\newcommand{\EventQFiveHundred}{17.60}
\newcommand{\CombinedQFiveHundred}{18.98}
\newcommand{\CombinedErrFiveHundred}{+7.8}
\newcommand{\EventMeanLossFiveHundred}{5.12}
\newcommand{\EventQFiveThousand}{17.04}
\newcommand{\CombinedQFiveThousand}{18.33}
\newcommand{\CombinedErrFiveThousand}{+7.6}
\newcommand{\EventMeanLossFiveThousand}{5.12}
\newcommand{\StaticQ}{16.81}
\newcommand{\ReplenishedQ}{18.26}
\newcommand{\ReplenishedEL}{5.15}
\newcommand{\EventCoverRec}{92}
\newcommand{\EventCoverEl}{76}
\newcommand{\EventYearOneLossRec}{32.2}
\newcommand{\EventYearOneLossRecHundred}{52.4}
\newcommand{\AnnualYearOneLossRec}{13.1}
\newcommand{\EventYearOneRetRec}{3.30}
\newcommand{\AnnualYearOneRetRec}{3.43}
\newcommand{\EventYearTenLossRec}{0.1}
\newcommand{\AnnualYearTenLossRec}{0.0}
\newcommand{\EventYearThirtyRetRec}{3.64}
\newcommand{\AnnualYearThirtyRetRec}{3.67}
\newcommand{\EventYearOneLossEl}{61.3}
\newcommand{\AnnualYearOneLossEl}{42.0}
\newcommand{\EventYearTenLossEl}{19.7}
\newcommand{\AnnualYearTenLossEl}{14.2}
\newcommand{\EventYearThirtyRetEl}{5.87}
\newcommand{\AnnualYearThirtyRetEl}{5.94}
\newcommand{\EventPaths}{200{,}000}

% Generated by scripts/tail_dependence.py; do not edit.
\newcommand{\TailTFiveRatio}{2.46}
\newcommand{\TailTFiveMean}{5.00}
\newcommand{\TailTFiveQNinetyNine}{28.85}
\newcommand{\TailTFiveQ}{41.42}
\newcommand{\TailTTenRatio}{1.90}
\newcommand{\TailTTenMean}{5.00}
\newcommand{\TailTTenQNinetyNine}{21.94}
\newcommand{\TailTTenQ}{31.91}
\newcommand{\TailTFiveRatioHigh}{1.72}
\newcommand{\TailTTenRatioHigh}{1.43}
\newcommand{\GaussQ}{16.81}
\newcommand{\GaussQNinetyNine}{12.68}

% Generated by scripts/extract_credit_risk.py; do not edit.
\newcommand{\MarginalLoanPDHigh}{0.4293}
\newcommand{\ScalarReconversionHigh}{5.100}
\newcommand{\MarginalLoanPDBasel}{0.4233}
\newcommand{\ScalarReconversionBasel}{5.030}
\newcommand{\HomogeneousLoanPD}{0.4207}


\title{Pricing Uncollateralized Microcredit on a Public Blockchain:\\Short Loans, Re-lending and a Locked First-Loss Reserve}
\author{Claude Code\thanks{Claude Code is an AI agent made by Anthropic, working with the
Microcredit Protocol Project; the research was directed by the project's human operator. See the
disclosure at the end of the paper. Sources, data and code:
\url{https://github.com/scottonchain/microcredit-theory}.}}
\date{Working paper, 5 October 2026}

\begin{document}
\maketitle

\begin{abstract}
\noindent A pooled lending protocol on a public blockchain lends stablecoins for 30 days without
collateral, and it funds a first-loss reserve from a share of the interest it collects. With no
default history, its premium and reserve must be set from scenarios. We calibrate a one-factor model
of portfolio default to such a pool and show how three features of the setting change the standard
calculation. Loans are short and rolled, so per-loan default rates must be converted to annual
probabilities: a default rate of 1\% per 30-day loan is an annual probability of 11.5\%. The pool
re-lends after write-offs, which raises the 99.9\% loss quantile by 4.5 to 23\% over the static model,
and we give the exact transform. A pool of 100 loans has a 99.9\% tail about a fifth above the
asymptotic one. A premium of 500 basis points over the funding rate carries annual default
probabilities up to 3.3\%, and probabilities of 3, 5 and 10\% need 600, 800 and 1,400 basis points.
For the reserve, we prove that lenders' cumulative return equals gross interest less fees, losses and
the reserve's final balance. The protocol never releases reserve funded by interest, so any reserve
share above the ratio of expected loss to the loan rate is a permanent transfer from lenders. At the
protocol's production settings, lenders' long-run expected return falls from \LongRunAtElFive\% to
\LongRunLockedFive\%, below the \Effr\% funding rate. In a simulation of pools of 100 to 5,000 loans
at the contract's own timing, with interest arriving loan by loan, the largest observed difference
in mean return from the annual recursion behind these numbers, from the fifth year on, is
\EventDiffMaxFiveHundredPlus\ percentage points in pools of 500 loans or more (\EventDiffMaxHundred\
points with 100 loans), and the first year, before the reserve has built, is riskier than the
recursion says. A
reserve share at or modestly above expected loss keeps most of lenders' return; whether further
protection is cheaper as external first-loss capital depends on the horizon and the hurdle rate,
which the companion paper examines. The quantiles are conditional on the Gaussian factor: a
Student-$t$ factor at the same correlation roughly doubles the 99.9\% quantile.
\end{abstract}

\medskip
\noindent\textbf{Keywords:} credit risk, Vasicek model, granularity adjustment, first-loss reserve,
microcredit, decentralized finance.\quad \textbf{JEL:} G21, G23, G32.

% ---------------------------------------------------------------------------------------------
\section{Introduction}

Lending protocols on public blockchains hold pooled deposits of stablecoins and lend them to
borrowers \citep{gudgeon2020}. The largest of them, Compound and Aave, require collateral worth more
than the loan \citep{leshner2019compound,aave2020}, so their credit risk is market risk on the collateral. This
paper studies a pool that lends without collateral, to borrowers who lack it, on credit lines granted
by an issuer and on backing from other accounts \citep{conservedcredit2026}. Its credit risk is the
classic kind: unsecured loans default, and lenders bear what the pool's buffers do not.

Three features of the pool set it apart from the bank portfolios that standard capital models were
built for. Loans run for 30 days and are rolled, so a per-loan default rate is not an annual
probability of default. The pool re-lends a slot as soon as its defaulted loan is written off, so it
can suffer more than one default per slot in a year. And it is small, a few hundred loans, so the
asymptotic single-factor model understates its tail. The pool also has a first-loss reserve, funded
from a share of every interest payment, whose rules are set by the protocol rather than by the
lenders: interest-funded reserve is never released to them.

The pool has no default history, so no parameter can be estimated. We treat the annual probability
of default (PD) and the asset correlation as scenarios and report every result conditional on them.
The contributions are the following.

\begin{itemize}
\item \emph{Horizons} (Proposition~\ref{prop:horizon}). Under a constant hazard, a per-loan default
  rate $p_T$ of a $T$-day loan is an annual PD of $1-(1-p_T)^{365/T}$, and defaults per exposure-year
  with replacement lie between the PD and the hazard. A ``97\% repayment rate'' on 30-day loans is an
  annual PD of 31\%.
\item \emph{Re-lending} (Proposition~\ref{prop:replenish}). The loss per exposure-year of a pool that
  replaces defaulted loans is an increasing function of the conditional PD, so its quantiles are the
  same function of the Vasicek quantiles. Re-lending raises the 99.9\% quantile by 9\% at an annual
  PD of 5\% and by 23\% at 20\%.
\item \emph{Pricing} (Section~\ref{sec:pricing}). Break-even premiums at the loan and pool level,
  with the cash drag of idle deposits, a protocol fee, re-lending and finite pools; the largest PD a
  given premium carries; and the effect of stake-secured principal.
\item \emph{A locked reserve} (Theorems~\ref{thm:identity} and~\ref{thm:longrun}). Lenders'
  cumulative return equals gross interest net of fees, less losses, less the reserve's final balance,
  under any release policy. When interest-funded reserve can never be released, lenders' long-run
  expected return is $u\bigl[a(1-f)-\max\{\EL,\,ar\}\bigr]$ for utilisation $u$, loan rate $a$, fee
  share $f$ and reserve share $r$. A reserve share above $\EL/a$ therefore costs lenders its excess
  every year, for good.
\item \emph{An event-timed check} (Section~\ref{sec:event}). A simulation of pools of 100 to 5,000
  loans at the contract's own timing, with interest arriving loan by loan and write-offs one late
  period after default, differs from the annual recursion behind the reserve results by at most
  \EventDiffMaxFiveHundredPlus\ percentage points of observed mean return from the fifth year on in
  pools of 500 loans or more (\EventDiffMaxHundred\ points with 100 loans; the largest differences
  observed, not bounds), shows that the first year is riskier than the recursion says, and finds
  the premium's combined re-lending and granularity approximation
  a few percent above the simulated 99.9\% loss quantile at the three sizes tested, on the
  conservative side. It is a check of a specified replenished model at three sizes, not an error
  bound over parameters.
\end{itemize}
The market equilibrium of the loan rate and the reserve share, when lenders and borrowers respond to
them, is the subject of a companion paper \citep{lendingequilibrium2026}, which restates
Theorem~\ref{thm:longrun} and builds on it; this paper is the canonical derivation of that result.

The last result has a direct consequence for the protocol. Its production settings pair an 800
basis-point premium with a reserve share of 65\%, chosen for a policy that released the reserve above
a cap. Under the lock the protocol actually enforces, lenders' long-run expected return at an annual
PD of 5\% is \LongRunLockedFive\%, below the funding rate of \Effr\%; with the reserve share at
expected loss it is \LongRunAtElFive\%.

Section~\ref{sec:pool} describes the pool. Section~\ref{sec:horizon} treats horizons,
Section~\ref{sec:losses} the loss distribution, Section~\ref{sec:pricing} pricing and
Section~\ref{sec:reserve} the reserve. Section~\ref{sec:limits} lists the limitations.

% ---------------------------------------------------------------------------------------------
\section{The Pool}\label{sec:pool}

Lenders deposit a stablecoin into one pool and receive shares; borrowers draw from the same pool.
Table~\ref{tab:pool} lists the terms that matter for credit risk. They are taken from the protocol's
contract \citep{microcredit2026contract}.

\begin{table}[ht]
  \centering\small
  \caption{Terms of the pool used in the calibration.}
  \label{tab:pool}
  \begin{tabular}{p{0.24\linewidth}p{0.68\linewidth}}
    \toprule
    Term & Value and treatment \\
    \midrule
    Loan rate & $a=$ funding rate $+$ premium, fixed at origination; the funding rate is the
      effective federal funds rate, 4.33\% in the calibration and in the deployments to date. Simple
      interest on the original principal, none if repaid within 24 hours. \\
    Term & 30 days by default (1 to 365 allowed), so $n=365/30\approx12.17$ loans per slot-year. \\
    Recognition & Cash basis: interest is booked when paid, so a write-off removes principal only. \\
    Default & A loan can be written off 30 days after its due date and can never be repaid
      afterwards. Loss given default (LGD) is 100\% on unsecured principal. \\
    Secured principal & Principal backed by staked stablecoins, which are slashed into the pool at
      default: LGD 0. \\
    Unsecured backing & Another account's credit, burned at default; it recovers no cash, so it does
      not change LGD. \\
    Utilisation & At most 90\% of lenders' assets lent; $u=85\%$ in the base case. \\
    Fee & A share $f$ of interest, at most 20\%; 0 by default. \\
    Reserve & A share $r$ of interest, at most 80\%, enters a first-loss reserve. The reserve is a
      junior claim inside the pool, excluded from the value of lenders' shares; it pays losses before
      them. Capital added by outsiders can be released to lenders; interest-funded reserve cannot. \\
    \bottomrule
  \end{tabular}
\end{table}

The last row is the one that matters most below. The protocol credits each borrower's share of the
reserve inflow to it as credit it can borrow against, and to keep that credit from being recaptured,
it never releases interest-funded reserve \citep{conservedcredit2026}. This paper takes the rule as
given and works out what it costs lenders.

Three configurations of the contract exist, and the paper names them. The \emph{production} settings
are the calibrated defaults of the production deployment script at commit \texttt{b725a85}: a premium
of 800 basis points and a reserve share of 65\%. The \emph{deployed} settings are those of the live
testnet deployment: a premium of 500 basis points and a share of 30\%. The development defaults differ
again. ``Deployed premium'' below means 500 basis points and ``production settings'' the first pair.
On 6 October 2026, after the pinned commit, the production script's default share became an interim
45\% (contract commit \texttt{1812e7d}), pending review of this paper and of the companion paper on the
market equilibrium \citep{lendingequilibrium2026}; Corollary~\ref{cor:production} describes the 65\% of
the pinned commit, and Remark~\ref{rem:interim} the interim value.

% ---------------------------------------------------------------------------------------------
\section{Default Horizons}\label{sec:horizon}

Write $p_T$ for the probability that a loan of $T$ days is never repaid, and PD, without a subscript,
for the annual probability that a borrower who keeps rolling such loans defaults within a year, the
horizon of the Basel framework \citep{bcbs2006}.

\begin{proposition}[Horizons]\label{prop:horizon}
Fix the systematic factor $Z$ for the year and suppose that, given $Z$, each of a borrower's
successive loans defaults with the same probability $p_T(Z)$, independently of earlier loans. Then
the conditional annual probability of default is
\[
  \mathrm{PD}(Z)=1-\bigl(1-p_T(Z)\bigr)^{n},\qquad n=365/T,
\]
the number of loans completed before default is geometric, and if every defaulter is replaced at
once the expected number of defaults per exposure-year given $Z$ is $n\,p_T(Z)$, which satisfies
$\mathrm{PD}(Z)\le n\,p_T(Z)\le\lambda(Z)$ with $\lambda(Z)=-\ln(1-\mathrm{PD}(Z))$ and increases with $n$.
\end{proposition}

\begin{proof}
Given $Z$, the borrower survives the year if it survives $n$ independent loans, and the count of
completed loans before the first default is geometric. With $x=\lambda/n$,
$n\,p_T=n(1-e^{-\lambda/n})=\lambda(1-e^{-x})/x$. The function $(1-e^{-x})/x$ decreases in $x>0$
from 1, so $n\,p_T$ increases in $n$, equals PD at $n=1$ and tends to $\lambda$ as $n\to\infty$.
\end{proof}

\begin{remark}[Conditional, marginal and homogeneous conversions]\label{rem:marginal}
The conversion holds given the factor. Unconditionally, $\mathrm{PD}=\E\bigl[1-(1-p_T(Z))^n\bigr]$,
which is not $1-(1-\E\,p_T(Z))^n$: the function is concave in $p_T$, so the scalar conversion applied to
the marginal per-loan rate overstates the annual PD. At an annual PD of 5\% and a correlation of 0.15 the
marginal per-loan rate is \MarginalLoanPDHigh\% and the scalar conversion returns \ScalarReconversionHigh\%;
at the Basel correlation the figures are \MarginalLoanPDBasel\% and \ScalarReconversionBasel\%, against
\HomogeneousLoanPD\% per loan for a homogeneous hazard. Table~\ref{tab:horizon} is the homogeneous
conversion, which is also the conditional one, and the pricing of Section~\ref{sec:pricing} uses
conditional quantities throughout. The exponential default time of hazard $\lambda$ and the fractional
cycle counts it allows are a continuous-hazard interpolation of the geometric count, adopted as a
convention; nothing below depends on it beyond the value of $n$.
\end{remark}

The assumption ignores two effects of opposite sign. Borrowers who have repaid are selected safer,
and a borrower may default strategically on its last loan, when future access is worth least.
Neither can be estimated without data. Table~\ref{tab:horizon} gives the conversion. The second
panel shows why per-loan rates mislead: 1\% per 30-day loan is an annual PD of 11.5\%, and a
``97\% repayment rate'' is 31\%. The conversion applies to a per-loan default rate with the same
denominator, loans of term $T$ each counted as defaulted when written off; an institution's reported
repayment rate has its own denominator, cohort and rollover rules, and the 31\% is an illustration of
scale, not an estimate for any institution. Default rates of this size are within the range reported
for microfinance institutions \citep{cull2009}, so the conversion matters for the scenarios below.

\begin{table}[ht]
  \centering\small
  \caption{Default horizons for 30-day loans under Proposition~\ref{prop:horizon}.}
  \label{tab:horizon}
  \begin{tabular}{rrrr}
    \toprule
    Annual PD & Per 30-day loan & Hazard $\lambda$ & Defaults per exposure-year, replaced \\
    \midrule
    \inputrows{data/table_horizon.tex}
    \bottomrule
  \end{tabular}
  \par\medskip
  \begin{tabular}{rrr}
    \toprule
    Per 30-day loan & Annual PD & Defaults per exposure-year, replaced \\
    \midrule
    \inputrows{data/table_loan_to_annual.tex}
    \bottomrule
  \end{tabular}
\end{table}

% ---------------------------------------------------------------------------------------------
\section{The Loss Distribution}\label{sec:losses}

\subsection{The one-factor model}

Borrower $i$ defaults within the year if $\sqrt{\rho}\,Z+\sqrt{1-\rho}\,\varepsilon_i<\Phi^{-1}(\mathrm{PD})$,
with $Z$ and the $\varepsilon_i$ independent standard normal \citep{merton1974,vasicek2002}.
Conditional on $Z$, defaults are independent with probability
\[
  p(Z)=\Phi\Bigl(\frac{\Phi^{-1}(\mathrm{PD})-\sqrt{\rho}\,Z}{\sqrt{1-\rho}}\Bigr),
\]
and in an infinitely granular pool the default rate converges to $D=p(Z)$, with loss rate
$L=\mathrm{LGD}\cdot D$. Its $q$-quantile is
$L_q=\mathrm{LGD}\cdot\Phi\bigl((\Phi^{-1}(\mathrm{PD})+\sqrt{\rho}\,\Phi^{-1}(q))/\sqrt{1-\rho}\bigr)$,
the expected loss is $\EL=\mathrm{PD}\cdot\mathrm{LGD}$, and the unexpected loss at 99.9\%,
$\UL=L_{0.999}-\EL$, is the Basel capital charge for retail exposures without its scaling factor
\citep{gordy2003,bcbs2005}. The expected shortfall and the stop-loss transform $\E[(D-k)^+]$ also
have closed forms in the bivariate normal distribution; Appendix~\ref{app:closed} gives them.

For the correlation we use the Basel formula for other retail exposures, which falls from 0.16 to
0.03 as PD rises \citep{bcbs2006}, and a grid from 0.02 to 0.30. The Basel value was calibrated on
bank retail books. A small pool whose borrowers share a community, a region or income tied to
crypto-assets is plausibly more correlated, which is what the upper end of the grid is for.

\subsection{Re-lending}

The Basel view is static: it follows the borrowers present at the start of the year. This pool
re-lends a slot as soon as its defaulted loan is written off.

\begin{proposition}[Re-lending]\label{prop:replenish}
Fix the factor for the year, and suppose each slot is lent in $n$ consecutive loans with a defaulter
replaced at once. Then the expected defaults per slot-year are $h\bigl(p(Z)\bigr)$ with
$h(p)=n\bigl(1-(1-p)^{1/n}\bigr)$, and the $q$-quantile of the loss per exposure-year in an infinitely
granular pool is $\mathrm{LGD}\cdot h(D_q)$, where $D_q$ is the Vasicek quantile.
\end{proposition}

\begin{proof}
Given $Z$, each loan defaults with the per-loan probability $1-(1-p(Z))^{1/n}$ of
Proposition~\ref{prop:horizon}, and a slot carries $n$ loans a year whatever happened to earlier
ones. The function $h$ is increasing on $[0,1]$, so it maps quantiles of $p(Z)$ to quantiles of
$h(p(Z))$.
\end{proof}

Re-lending changes little at the mean, 5.15\% against 5.00\% at an annual PD of 5\%, and more in the
tail: the 99.9\% quantile rises by 4.5\% at PD 1\%, 8.7\% at 5\% and 23\% at 20\%
(Table~\ref{tab:losses}). A revolving Monte Carlo with 500 slots, a one-term lag before write-off and
the factor fixed for the year reproduces both the static and the replenished values. We use the
static distribution in the tables, because it is the standard and citable object, and the
replenished one in the pricing recommendation, because it is more conservative.

\begin{table}[ht]
  \centering\small
  \caption{Loss rate per unit of exposure at the Basel correlation (LGD 100\%). The last two columns
  are for a pool that re-lends after write-offs (Proposition~\ref{prop:replenish}).}
  \label{tab:losses}
  \begin{tabular}{rrrrrrrrr}
    \toprule
    PD & $\rho$ & EL & s.d. & $L_{0.99}$ & $L_{0.999}$ & ES$_{0.999}$ & EL, re-lent & $L_{0.999}$, re-lent \\
    \midrule
    \inputrows{data/table_losses.tex}
    \bottomrule
  \end{tabular}
\end{table}

\begin{table}[ht]
  \centering\small
  \caption{The 99.9\% loss quantile on the correlation grid.}
  \label{tab:rhogrid}
  \begin{tabular}{rrrrrrrr}
    \toprule
    PD & Basel & 0.02 & 0.05 & 0.10 & 0.15 & 0.20 & 0.30 \\
    \midrule
    \inputrows{data/table_rho_grid.tex}
    \bottomrule
  \end{tabular}
\end{table}

\begin{figure}[ht]
  \centering
  % Figure: density of the annual loss rate at PD 5% (data/loss_density.csv; density per percentage point).
\begin{tikzpicture}
  \begin{axis}[
      width=0.72\linewidth, height=5.2cm,
      xmin=0, xmax=25, ymin=0,
      xlabel={annual loss rate (\%)}, ylabel={density per point},
      tick label style={font=\footnotesize}, label style={font=\small},
      legend style={font=\footnotesize, draw=none, at={(0.97,0.95)}, anchor=north east},
      scaled y ticks=false, yticklabel style={/pgf/number format/fixed, /pgf/number format/precision=2},
    ]
    \addplot[thick, blue!65!black] table[x=loss_pct, y=basel, col sep=comma] {data/loss_density.csv};
    \addlegendentry{$\rho=0.053$ (Basel)}
    \addplot[thick, dashed, red!70!black] table[x=loss_pct, y=rho015, col sep=comma] {data/loss_density.csv};
    \addlegendentry{$\rho=0.15$}
  \end{axis}
\end{tikzpicture}

  \caption{Density of the annual loss rate at an annual PD of 5\%, at the Basel correlation (0.053)
  and at 0.15. The mean is 5\% in both; the higher correlation moves mass into the tail.}
  \label{fig:density}
\end{figure}

\subsection{Finite pools}

A pool of 10,000 units lent in loans of 100 holds about 100 loans. Drawing $Z$ and then the number of
defaults as Binomial$(N,p(Z))$ is exact for equal exposures, and a brute-force simulation of every
borrower's latent variable confirms it. The analytic comparison is the first-order granularity
adjustment of \citet{martin2002} and \citet{gordy2004}; see also \citet{gordy2013}. With $10^6$
scenarios per cell, quantile confidence intervals come from order statistics.
Table~\ref{tab:granularity} shows that at 100 loans the 99.9\% quantile is about a fifth above the
asymptotic value, that the adjustment captures most of the gap, and that at 5,000 loans the gap is
negligible. With 100 loans, losses come in steps of 1\%, so a confidence interval can collapse onto
one point of the lattice.

\begin{table}[ht]
  \centering\small
  \caption{The 99.9\% loss quantile of a finite pool at the Basel correlation: asymptotic,
  asymptotic with the granularity adjustment, and Monte Carlo with a 95\% interval ($10^6$ scenarios).}
  \label{tab:granularity}
  \begin{tabular}{rrrrl}
    \toprule
    PD & Loans $N$ & Asymptotic & With adjustment & Monte Carlo \\
    \midrule
    \inputrows{data/table_granularity.tex}
    \bottomrule
  \end{tabular}
\end{table}

\subsection{Beyond the Gaussian factor}\label{sec:tail}

The Gaussian factor is a choice, not a bound. In the $t$-copula alternative the latent variable is
$\sqrt{W}\,(\sqrt{\rho}\,Z+\sqrt{1-\rho}\,\varepsilon_i)$ with $W=\nu/\chi^2_\nu$ independent of
$(Z,\varepsilon)$, and the threshold is the $t_\nu$ quantile of PD, so the marginal PD and the linear
correlation of the latent variables are unchanged while defaults share a tail factor
\citep{frey2003,mcneil2015}. Conditional on $(W,Z)$ defaults are independent with probability
$\Phi\bigl((t_\nu^{-1}(\mathrm{PD})/\sqrt W-\sqrt\rho\,Z)/\sqrt{1-\rho}\bigr)$, which is the loss rate
of the infinitely granular pool. Table~\ref{tab:tail} gives its 99\% and 99.9\% quantiles by Monte Carlo
(2,000,000 draws) against the Gaussian ones. At an annual PD of 5\% and the Basel correlation, 10
degrees of freedom raise the 99.9\% quantile from \GaussQ\% to \TailTTenQ\%, a factor of
\TailTTenRatio, and 5 degrees of freedom to \TailTFiveQ\%, a factor of \TailTFiveRatio; the mean is
\TailTFiveMean\% in each, as it must be. The correlation sweep of Table~\ref{tab:rhogrid} does not reach
these values. The paper's quantiles and premiums are therefore conditional on the Gaussian factor;
Section~\ref{sec:limits} says so, and a pool with reason to expect a shared tail should price with a
dependence assumption of its own.

\begin{table}[ht]
  \centering\small
  \caption{Loss quantiles of the infinitely granular pool (LGD 100\%) with a Gaussian factor and with
  a Student-$t$ factor of 5 and 10 degrees of freedom at the same linear correlation; the marginal PD
  is the same in each column.}
  \label{tab:tail}
  \begin{tabular}{rlrrrrrr}
    \toprule
    & & \multicolumn{2}{c}{Gaussian} & \multicolumn{2}{c}{$t$, $\nu=5$} & \multicolumn{2}{c}{$t$, $\nu=10$} \\
    \cmidrule(lr){3-4}\cmidrule(lr){5-6}\cmidrule(lr){7-8}
    PD & $\rho$ & $L_{0.99}$ & $L_{0.999}$ & $L_{0.99}$ & $L_{0.999}$ & $L_{0.99}$ & $L_{0.999}$ \\
    \midrule
    \inputrows{data/table_tail.tex}
    \bottomrule
  \end{tabular}
\end{table}

% ---------------------------------------------------------------------------------------------
\section{Pricing}\label{sec:pricing}

\subsection{Break-even premiums}

The loan rate is the funding rate plus a premium $s$. Lenders' opportunity cost is the funding
rate, so it passes through, and the loan-level break-even premium is
\begin{equation}\label{eq:sloan}
  s^*=\EL+h\cdot\UL,
\end{equation}
where $h$ is the required return in excess of the funding rate on risk capital $\UL$. In this pool
the lenders are the capital, since there is no equity tranche. The pool adds two costs. Lenders earn
nothing on the share $1-u$ of deposits held idle for withdrawals, and they pay the fee $f$ on interest.

\begin{proposition}[Pool-level premium]\label{prop:pool}
Per unit of deposits, lenders earn $u\,a(1-f)-u\,\EL$ in expectation and require the funding rate
$e$ plus $h\,u\,\UL$ on the capital at risk. With $a=e+s$, the break-even premium is
\[
  s_{\mathrm{pool}}=\frac{e/u+\EL+h\,\UL}{1-f}-e.
\]
\end{proposition}

\begin{proof}
Solve $u(e+s)(1-f)-u\,\EL=e+h\,u\,\UL$ for $s$.
\end{proof}

At $u=85\%$ the cash drag alone adds $e(1/u-1)=76$ basis points. The formula is first order: it
ignores interest lost on a defaulting loan, at most one term plus the late period, and compounding.
Table~\ref{tab:premium} reports the premiums at a 15\% hurdle, Table~\ref{tab:maxpd} the largest PD
that the deployed 500 basis points carries, and Figure~\ref{fig:premium} the premium against PD.

\begin{table}[ht]
  \centering\small
  \caption{Break-even premium in basis points at a hurdle of 15\%: loan level at four correlations,
  then at the Basel correlation re-lent, with 100 loans, and at the pool level ($u=85\%$, without and
  with a 10\% fee).}
  \label{tab:premium}
  \begin{tabular}{rrrrrrrrr}
    \toprule
    & \multicolumn{4}{c}{Loan level, $\rho$} & \multicolumn{2}{c}{Basel $\rho$} & \multicolumn{2}{c}{Pool level} \\
    \cmidrule(lr){2-5}\cmidrule(lr){6-7}\cmidrule(lr){8-9}
    PD & Basel & 0.05 & 0.15 & 0.30 & re-lent & $N=100$ & $f=0$ & $f=10\%$ \\
    \midrule
    \inputrows{data/table_premium.tex}
    \bottomrule
  \end{tabular}
\end{table}

\begin{table}[ht]
  \centering\small
  \caption{Largest annual PD whose break-even premium is at most 500 basis points, by hurdle; the
  last column is the pool level at $u=85\%$ and $h=15\%$.}
  \label{tab:maxpd}
  \begin{tabular}{lrrrrr}
    \toprule
    $\rho$ & $h=0$ & $h=10\%$ & $h=15\%$ & $h=25\%$ & Pool, $h=15\%$ \\
    \midrule
    \inputrows{data/table_maxpd.tex}
    \bottomrule
  \end{tabular}
\end{table}

\begin{figure}[ht]
  \centering
  % Figure: loan-level break-even premium against annual PD (data/premium_curve.csv).
\begin{tikzpicture}
  \begin{axis}[
      width=0.72\linewidth, height=5.6cm,
      xmode=log, xmin=0.5, xmax=20,
      xtick={0.5,1,2,3,5,10,20}, xticklabels={0.5,1,2,3,5,10,20},
      ymin=0, ymax=3200,
      xlabel={annual PD (\%)}, ylabel={break-even premium (bps)},
      tick label style={font=\footnotesize}, label style={font=\small},
      legend style={font=\footnotesize, draw=none, at={(0.03,0.97)}, anchor=north west},
      scaled y ticks=false, yticklabel style={/pgf/number format/fixed, /pgf/number format/1000 sep={,}},
    ]
    \addplot[thick, blue!65!black] table[x=pd_pct, y=basel, col sep=comma] {data/premium_curve.csv};
    \addlegendentry{Basel $\rho$}
    \addplot[thick, green!45!black] table[x=pd_pct, y=rho005, col sep=comma] {data/premium_curve.csv};
    \addlegendentry{$\rho=0.05$}
    \addplot[thick, orange!80!black] table[x=pd_pct, y=rho015, col sep=comma] {data/premium_curve.csv};
    \addlegendentry{$\rho=0.15$}
    \addplot[thick, red!70!black] table[x=pd_pct, y=rho030, col sep=comma] {data/premium_curve.csv};
    \addlegendentry{$\rho=0.30$}
    \addplot[dashed, gray, domain=0.5:20] {500};
  \end{axis}
\end{tikzpicture}

  \caption{Loan-level break-even premium against annual PD at a 15\% hurdle. The dashed line is the
  deployed premium of 500 basis points.}
  \label{fig:premium}
\end{figure}

A premium of 500 basis points covers expected loss up to an annual PD of 5\% and nothing more. With
a 15\% hurdle it carries PD up to 3.30\% at the loan level, or 0.28\% per 30-day loan, and 2.60\% once
the cash drag is counted. At the higher correlation of 0.15 the two limits are 2.39\% and 1.94\%.
Table~\ref{tab:recommendation} gives the premiums we recommend for the protocol at three scenario
PDs: the pool-level premium, re-lent and adjusted for a pool of 500 loans, rounded up to 50 basis
points. The contract has one premium for every loan, so it must be set for the riskiest PD it
admits, or admission must keep the book within the PD it was priced for.

\begin{table}[ht]
  \centering\small
  \caption{Premiums at the Basel correlation and a 15\% hurdle, in basis points, and the
  recommendation.}
  \label{tab:recommendation}
  \begin{tabular}{rrrrrrr}
    \toprule
    PD & Per loan & Loan level & Pool level & Pool, re-lent, $N=500$ & Recommended & Loan rate \\
    \midrule
    \inputrows{data/table_recommendation.tex}
    \bottomrule
  \end{tabular}
\end{table}

\subsection{Secured principal and backing}

If a share $\sigma$ of every loan's principal is secured by stake, LGD is $1-\sigma$. The loss
quantile is linear in LGD, so $\EL$, $\UL$ and $s^*$ all scale by $1-\sigma$, and with different
shares per loan the asymptotic model adds them loan by loan \citep{gordy2003}. At the Basel
correlation and a 15\% hurdle, 500 basis points breaks even with a secured share of 26\% at PD 5\% and
58\% at PD 10\%.

Unsecured backing recovers nothing, so it cannot lower LGD. If it has an effect, it is on PD: a
backer who loses credit when the borrower defaults has reason to screen and monitor
\citep{stiglitz1990peer,ghatak1999joint}. That effect is a hypothesis to estimate. A field experiment
found no rise in default when a lender removed group liability \citep{gine2014group}, so it may be
small. A single loan rate for backed and unbacked loans overprices backed loans if the effect is
real, which invites adverse selection among the unbacked.

% ---------------------------------------------------------------------------------------------
\section{The First-Loss Reserve}\label{sec:reserve}

\subsection{Sizing}

A share $r$ of interest enters the reserve. Per unit of deposits its yearly inflow is $u\,a\,r$ and
the losses it faces are $u\,L$, so utilisation cancels: a reserve that covers a loss rate $T$ needs
$r=T/a$, which is feasible from interest only if $T/a\le1-f$. The reserve pays from the first unit of
loss, so a reserve that covers $L_q-\EL$ keeps lenders' loss at most $u\,\EL$, the loss already
priced, with probability $q$. Covering expected loss takes $r_{\EL}=\EL/a$, which at the recommended
loan rates and re-lent losses is \RElThree\%, \RElFive\% and \RElTen\% at annual PDs of 3, 5 and 10\%.
The calibration that preceded this paper recommended shares of \RRecThree\%, \RRecFive\% and
\RRecTen\%: expected loss plus a third of the 99\% unexpected loss, with the balance capped at
$u(L_{0.99}-\EL)$ and the surplus released to lenders. The protocol adopted the 65\% share for its
production settings. It does not release interest-funded reserve.

\subsection{An accounting identity}

\begin{theorem}[Lenders' cumulative return]\label{thm:identity}
Suppose the reserve is funded only from interest, pays only losses and releases to lenders, and starts
empty. Then over any $T$ years and for every path of losses, lenders' cumulative return per unit of
deposits is
\[
  \sum_{t=1}^{T}\mathrm{APY}_t=T\,u\,a(1-f)-\sum_{t=1}^{T}u\,L_t-R_T,
\]
where $L_t$ is the loss rate in year $t$ and $R_T$ the reserve's balance at the end.
\end{theorem}

\begin{proof}
In year $t$ lenders receive $u\,a(1-f-r)$, bear the loss the reserve does not cover, $\ell_t$, and
receive any release $x_t$. The reserve receives $u\,a\,r$, pays $u\,L_t-\ell_t$ and releases $x_t$,
so $R_T=T\,u\,a\,r-\sum_t(u\,L_t-\ell_t)-\sum_tx_t$. Then
$\sum_t\mathrm{APY}_t=T\,u\,a(1-f-r)-\sum_t\ell_t+\sum_tx_t$, and substituting
$\sum_t\ell_t-\sum_tx_t=\sum_tu\,L_t+R_T-T\,u\,a\,r$ gives the claim.
\end{proof}

A reserve funded from lenders' own interest thus changes their return only through its final balance.
It moves losses in time and smooths the share price, but whatever it still holds is lenders' money
they will not receive. Under a policy that caps the balance and releases the rest, $R_T$ is bounded,
and lenders' long-run return is that of no reserve. Under the protocol's rule it is not.

\subsection{The locked reserve}

\begin{theorem}[Long-run return under a lock]\label{thm:longrun}
Suppose the loss rates $L_t$ are independent across years, take values in $[0,\bar L]$ for some
finite $\bar L$ (for the re-lent loss per exposure-year of Proposition~\ref{prop:replenish},
$\bar L=n$, since a slot can default at most once per loan), and have a common mean $\EL$, and the
reserve is never released. Then lenders' average annual return converges almost surely:
\[
  \frac{1}{T}\sum_{t=1}^T\mathrm{APY}_t\;\longrightarrow\;u\bigl[a(1-f)-\max\{\EL,\ a\,r\}\bigr].
\]
\end{theorem}

\begin{proof}
With $X_t=u\,a\,r-u\,L_t$, the balance follows the Lindley recursion $R_t=\max\{R_{t-1}+X_t,0\}$
\citep{lindley1952}, with $R_0=0$. Let $S_t=\sum_{s\le t}X_s$ and $\mu=\E X_t=u(a\,r-\EL)$. Unrolling the
recursion, $R_T=\max\{S_T,\ \max_{1\le j\le T}(S_T-S_j)\}=S_T-\min_{0\le j\le T}S_j$ with $S_0=0$. The
$X_t$ are independent and bounded, so $\sum_t\mathrm{Var}(X_t)/t^2<\infty$ and Kolmogorov's strong law gives
$S_n/n\to\mu$ almost surely \citep[Theorem~22.7]{billingsley1995}; the loss rates themselves give
$T^{-1}\sum_tL_t\to\EL$ in the same way. If $\mu>0$, then almost surely $S_n\to\infty$, so $\min_jS_j$
is finite and $R_T/T\to\mu$. If $\mu\le0$, fix $\epsilon>0$; almost surely there is $N$ with
$|S_n-n\mu|\le\epsilon n$ for all $n\ge N$. For $T\ge j\ge N$,
$S_T-S_j\le(T-j)\mu+\epsilon(T+j)\le2\epsilon T$; for $j<N\le T$, $S_T-S_j\le\epsilon T+T\mu+|S_j|\le
\epsilon T+C$ with $C=\max_{j<N}|S_j|$, which is finite. So $0\le R_T\le2\epsilon T+C$ for $T\ge N$,
$\limsup_T R_T/T\le2\epsilon$, and since $\epsilon$ is arbitrary, $R_T/T\to0$. In both cases
$R_T/T\to\max\{\mu,0\}$.
Dividing Theorem~\ref{thm:identity} by $T$ gives the limit $u\,a(1-f)-u\,\EL-\max\{\mu,0\}$, which is the
claim. The argument uses only the strong law for the partial sums, so it covers any sequence of losses
for which $S_n/n\to\mu$ and $T^{-1}\sum L_t\to\EL$ almost surely, such as a stationary ergodic one.
\end{proof}

A reserve share below $\EL/a$ leaves lenders' long-run return unchanged, because the reserve is
emptied by losses as fast as it fills. Above $\EL/a$, every additional unit of share costs lenders
$u\,a$ per year, for good: the excess accumulates in a balance they cannot receive. In the protocol
that balance is not idle in purpose. It backs the credit that the protocol grants borrowers for the
interest they have paid, and it absorbs defaults on that credit. The reserve share is therefore also
the rate at which lenders' yield is converted into borrowers' credit, a choice of policy rather than
of risk management.

Figure~\ref{fig:years} and Table~\ref{tab:policy} simulate both policies at an annual PD of 5\%, with
independent years, re-lent losses and 200,000 paths per case. The simulation is an annual recursion,
and its assumptions should be stated: deposits are constant, so losses are replenished by new deposits
and the exposure $u$ per unit of deposits does not move; each year's inflow $u\,a\,r$ arrives before
that year's loss $u\,L_t$ is paid; the loss rate is the conditional, infinitely granular, replenished
rate of Proposition~\ref{prop:replenish}, so a defaulted slot is replaced at once; and years are
independent. Section~\ref{sec:event} relaxes the timing, the instant replacement and the granularity
in a finite pool and reports what changes. Theorem~\ref{thm:identity} holds on every path to rounding
error. In the recursion, at the production share of 65\%, the two policies give the same first year,
3.43\%, because the reserve is filling. By year 5 the policy with release pays lenders 5.93\% and the
locked one 3.66\%, and after 30 years the locked reserve holds 73\% of deposits. At the expected-loss
share of 42\%, the two policies stay within 0.2 percentage points of each other, and from year 5 on
lenders bear some loss in 8 to 20\% of years.

\begin{table}[ht]
  \centering\small
  \caption{Lenders' mean annual return and the probability that they bear any loss in the year, by
  year, at an annual PD of 5\% (Basel correlation, loan rate 12.33\%, $u=85\%$, $f=0$); the last column
  is the mean reserve after 30 years as a share of deposits.}
  \label{tab:policy}
  \begin{tabular}{lrrrrrrrrr}
    \toprule
    & \multicolumn{2}{c}{Year 1} & \multicolumn{2}{c}{Year 5} & \multicolumn{2}{c}{Year 10} & \multicolumn{2}{c}{Year 30} & Reserve \\
    \cmidrule(lr){2-3}\cmidrule(lr){4-5}\cmidrule(lr){6-7}\cmidrule(lr){8-9}
    Reserve share and policy & Return & Loss & Return & Loss & Return & Loss & Return & Loss & year 30 \\
    \midrule
    \inputrows{data/table_reserve_policy.tex}
    \bottomrule
  \end{tabular}
\end{table}

\begin{figure}[ht]
  \centering
  % Figure: lenders' return and loss probability by year at PD 5% (data/apy_by_year.csv).
\begin{tikzpicture}
  \begin{groupplot}[
      group style={group size=2 by 1, horizontal sep=1.6cm},
      width=0.47\linewidth, height=5.2cm,
      xmin=1, xmax=30, xlabel={year},
      tick label style={font=\footnotesize}, label style={font=\small},
      legend style={font=\scriptsize, draw=none, fill=white, fill opacity=0.85, text opacity=1},
    ]
    \nextgroupplot[ylabel={mean return (\%)}, ymin=3, ymax=6.6, legend style={at={(0.97,0.03)}, anchor=south east}]
    \addplot[thick, blue!65!black] table[x=year, y=el_rel_apy, col sep=comma] {data/apy_by_year.csv};
    \addlegendentry{42\%, released}
    \addplot[thick, dashed, blue!65!black] table[x=year, y=el_lock_apy, col sep=comma] {data/apy_by_year.csv};
    \addlegendentry{42\%, locked}
    \addplot[thick, red!70!black] table[x=year, y=rec_rel_apy, col sep=comma] {data/apy_by_year.csv};
    \addlegendentry{65\%, released}
    \addplot[thick, dashed, red!70!black] table[x=year, y=rec_lock_apy, col sep=comma] {data/apy_by_year.csv};
    \addlegendentry{65\%, locked}
    \addplot[dotted, gray, domain=1:30] {4.33};

    \nextgroupplot[ylabel={P(lenders bear a loss) (\%)}, ymin=0, ymax=45]
    \addplot[thick, blue!65!black] table[x=year, y=el_rel_p, col sep=comma] {data/apy_by_year.csv};
    \addplot[thick, dashed, blue!65!black] table[x=year, y=el_lock_p, col sep=comma] {data/apy_by_year.csv};
    \addplot[thick, red!70!black] table[x=year, y=rec_rel_p, col sep=comma] {data/apy_by_year.csv};
    \addplot[thick, dashed, red!70!black] table[x=year, y=rec_lock_p, col sep=comma] {data/apy_by_year.csv};
  \end{groupplot}
\end{tikzpicture}

  \caption{Lenders' mean annual return (left) and the probability that they bear any loss in the
  year (right), at an annual PD of 5\%, for reserve shares of 42\% (expected loss) and 65\% (the
  production setting), each with release above a cap and locked.}
  \label{fig:years}
\end{figure}

Figure~\ref{fig:sweep} shows the trade-off the lock creates. The long-run return is flat in the
reserve share up to $\EL/a$ and falls linearly beyond it, as Theorem~\ref{thm:longrun} states, while
the probability of a loss in year 10 falls steeply just above $\EL/a$. At an annual PD of 5\%, a share
of 50\%, about 1.2 times $\EL/a$, costs lenders 0.9 percentage points of long-run return and lowers the
probability of a loss in year 10 from 14\% at the expected-loss share to 2.5\%. The production share of
65\% costs 2.4 points and lowers it to 0.03\%.

\begin{figure}[ht]
  \centering
  % Figure: locked reserve, long-run return and loss probability against the reserve share
% (data/reserve_sweep_curves.csv).
\begin{tikzpicture}
  \begin{groupplot}[
      group style={group size=2 by 1, horizontal sep=1.6cm},
      width=0.47\linewidth, height=5.2cm,
      xmin=0.2, xmax=0.8, xlabel={reserve share $r$},
      tick label style={font=\footnotesize}, label style={font=\small},
      legend style={font=\scriptsize, draw=none, fill=white, fill opacity=0.85, text opacity=1},
    ]
    \nextgroupplot[ylabel={mean return, year 30 (\%)}, ymin=0, ymax=7.5, legend style={at={(0.03,0.03)}, anchor=south west}]
    \addplot[thick, green!45!black] table[x=r, y=apy3, col sep=comma] {data/reserve_sweep_curves.csv};
    \addlegendentry{PD 3\%}
    \addplot[thick, blue!65!black] table[x=r, y=apy5, col sep=comma] {data/reserve_sweep_curves.csv};
    \addlegendentry{PD 5\%}
    \addplot[thick, red!70!black] table[x=r, y=apy10, col sep=comma] {data/reserve_sweep_curves.csv};
    \addlegendentry{PD 10\%}
    \addplot[only marks, mark=o, mark size=1.2pt, green!45!black, forget plot] table[x=r, y=f3, col sep=comma] {data/reserve_sweep_curves.csv};
    \addplot[only marks, mark=o, mark size=1.2pt, blue!65!black, forget plot] table[x=r, y=f5, col sep=comma] {data/reserve_sweep_curves.csv};
    \addplot[only marks, mark=o, mark size=1.2pt, red!70!black, forget plot] table[x=r, y=f10, col sep=comma] {data/reserve_sweep_curves.csv};
    \addplot[dotted, gray, domain=0.2:0.8] {4.33};

    \nextgroupplot[ylabel={P(lenders bear a loss), year 10 (\%)}, ymin=0, ymax=100]
    \addplot[thick, green!45!black] table[x=r, y=p3, col sep=comma] {data/reserve_sweep_curves.csv};
    \addplot[thick, blue!65!black] table[x=r, y=p5, col sep=comma] {data/reserve_sweep_curves.csv};
    \addplot[thick, red!70!black] table[x=r, y=p10, col sep=comma] {data/reserve_sweep_curves.csv};
  \end{groupplot}
\end{tikzpicture}

  \caption{Locked reserve: lenders' mean return in year 30 (left; markers are the limit of
  Theorem~\ref{thm:longrun}) and the probability of any lender loss in year 10 (right), against the
  reserve share, at annual PDs of 3, 5 and 10\% and the recommended loan rates.}
  \label{fig:sweep}
\end{figure}

\begin{corollary}[The production settings]\label{cor:production}
At the production settings (premium 800 basis points, reserve share 65\%, $u=85\%$, $f=0$) and an
annual PD of 5\%, lenders' long-run expected return under the lock is \LongRunLockedFive\%, against
\LongRunAtElFive\% with the reserve share at expected loss and a funding rate of \Effr\%. The same
comparison at PDs of 3 and 10\%, with the premiums and shares recommended for them, gives
\LongRunLockedThree\% against \LongRunAtElThree\% and \LongRunLockedTen\% against \LongRunAtElTen\%.
\end{corollary}

\begin{proof}
Theorem~\ref{thm:longrun} with the values of Table~\ref{tab:recommendation} and the re-lent expected
losses of Table~\ref{tab:losses}.
\end{proof}

\begin{remark}[The interim default]\label{rem:interim}
On 6 October 2026 the production script's default share became an interim 45\% (contract commit
\texttt{1812e7d}), pending review; the pinned commit the corollary describes had 65\%, and the live
testnet deployment keeps 30\%. At 45\% and an annual PD of 5\%, lenders' long-run return under the lock
is \LongRunFortyFiveFive\%, above the funding rate and \LongRunAtElFive\% at expected loss, because 45\%
is a few points above the expected-loss share of \RElFive\%.
\end{remark}

Under the lock at the pinned settings, a lender earns less than the funding rate in the long run at
every scenario PD, so the pool would not attract deposits that have the funding rate as their
alternative. Three remedies follow from the analysis. The reserve share can be set at or modestly above
$\EL/a$, which keeps lenders' return and accepts that they bear some losses in bad years. Protection
beyond that can come from external first-loss capital, which the protocol accepts from anyone, absorbs
losses before lenders, and, unlike interest-funded reserve, can release to lenders once it exceeds what
the pool needs. Or the premium can be raised to pay for the lock, which by Proposition~\ref{prop:pool}
means adding the excess share $\max\{0,r-\EL/a\}$ of the loan rate to the costs the premium must cover.
Which share and rate clear a market in which lenders and borrowers both respond is the question of the
companion paper \citep{lendingequilibrium2026}, which finds a volume-maximising share a few points above
$\EL/a$ under its scenarios.

\subsection{An event-timed check in a finite pool}\label{sec:event}

The recursion's assumptions matter most while the reserve forms, which is where its protection claims
are strongest. To test them we simulate the pool at the contract's own timing
\citep{microcredit2026contract}: $N$ slots of equal principal are lent in back-to-back loans of $365/12$
days, so that each year holds twelve terms; at the end of a term a loan repays its principal with simple
interest, of which the share $r$ goes to the reserve and the rest to lenders, or defaults; a defaulted
loan is written off one late period later, when the reserve pays its principal first and lenders bear
the rest, and only then is the slot re-lent; the systematic factor is fixed for the year, and the
defaults among the active slots in a term are binomial with the conditional per-loan probability of
Proposition~\ref{prop:horizon}. Deposits are constant at $N/u$, so exposure is replenished as in the
recursion, and the write-offs of a term are paid before that term's interest arrives. Two
simplifications remain: all loans start and end at term boundaries (the contract's are staggered),
and the term is $365/12$ days rather than the contract's 30, so that calendar years and terms align;
there are no partial repayments, impairments or withdrawals. Interest is what the active, repaying
slots pay, not the recursion's deterministic $u\,a$ a year, so the model differs from the recursion
in its inflows as well as their timing. Table~\ref{tab:event} reports \EventPaths\ paths for pools of
100, 500 and 5,000 loans at the two shares of Table~\ref{tab:policy}.

Three things change. First, the first year is worse than the recursion says: at the
production share, lenders bear some loss in \EventYearOneLossRec\% of paths with 500 loans and
\EventYearOneLossRecHundred\% with 100, against \AnnualYearOneLossRec\% in the recursion, because the
first defaults arrive before much interest has and because a small pool's losses are lumpy. Even so, the
reserve pays \EventCoverRec\% of first-year write-offs with 500 loans (\EventCoverEl\% at the
expected-loss share). Second, the gap narrows with time but does not close everywhere: at the
production share the probability of a lender loss in year 10 is \EventYearTenLossRec\% against
\AnnualYearTenLossRec\%, and from year 5 on the mean returns differ by at most
\EventDiffMaxRecFiveHundred\ points with 500 loans (\EventDiffMaxRecHundred\ with 100 and
\EventDiffMaxRecFiveThousand\ with 5,000); at the expected-loss share, where the reserve hovers near
empty, the event-timed loss probability stays above the recursion's (\EventYearTenLossEl\% against
\AnnualYearTenLossEl\% in year 10), because within a year losses can arrive before the interest that
would have covered them, and the mean returns differ by at most \EventDiffMaxElFiveHundred\ points
with 500 loans (\EventDiffMaxElHundred\ with 100 and \EventDiffMaxElFiveThousand\ with 5,000). These
are the largest differences observed between simulated means, rounded to two decimals, not bounds,
and they fall in years \EventDiffYearMin\ to \EventDiffYearMax. Third, the write-off
delay idles a defaulted slot for one term, which lowers the mean loss rate slightly, from
\ReplenishedEL\% to \EventMeanLossFiveHundred\%, and lowers the 99.9\% annual loss quantile below the
replenished one. Table~\ref{tab:eventq} compares the event-timed quantile with the recommendation's
combined approximation, the replenished quantile plus the static granularity adjustment: the
approximation overstates the quantile by \CombinedErrFiveHundred\% at 500 loans and
\CombinedErrHundred\% at 100, so it errs on the conservative side, by less than the differences between
correlation assumptions in Table~\ref{tab:rhogrid}. The long-run return moves least: the event-timed
mean return in year 30 is \EventYearThirtyRetRec\% against \AnnualYearThirtyRetRec\% at the production
share and \EventYearThirtyRetEl\% against \AnnualYearThirtyRetEl\% at expected loss, the differences
coming from the interest that late and defaulted slots do not pay. Theorem~\ref{thm:longrun} is stated
for the recursion's deterministic interest and does not by itself give the limit of this process; a
limiting statement for it would need the pathwise accounting of Theorem~\ref{thm:identity} with the
actual interest and loss inflows, which we do not pursue here. The agreement is a Monte Carlo
comparison at three pool sizes and two shares, under this model's own simplifications, not an error
bound over parameters.

The check keeps deposits constant. A pool whose deposits fall with losses and withdrawals, so that
exposure and the reserve's inflow shrink after a bad year, is the remaining difference between the
recursion and the protocol, and it is not tested here; the recursion's results are conditional on it.

\begin{table}[ht]
  \centering\small
  \caption{Event-timed simulation against the annual recursion at an annual PD of 5\% (Basel
  correlation, loan rate 12.33\%, $u=85\%$, $f=0$, locked reserve): lenders' mean annual return and the
  probability that they bear any loss in the year, for years 1, 5, 10 and 30, and the share of
  first-year write-offs that the reserve paid.}
  \label{tab:event}
  \begin{tabular}{lrrrrrrrrrr}
    \toprule
    & & \multicolumn{2}{c}{Year 1} & \multicolumn{2}{c}{Year 5} & \multicolumn{2}{c}{Year 10} & \multicolumn{2}{c}{Year 30} & Reserve \\
    \cmidrule(lr){3-4}\cmidrule(lr){5-6}\cmidrule(lr){7-8}\cmidrule(lr){9-10}
    Share and model & $N$ & Return & Loss & Return & Loss & Return & Loss & Return & Loss & paid, yr 1 \\
    \midrule
    \inputrows{data/table_event_timed.tex}
    \bottomrule
  \end{tabular}
\end{table}

\begin{table}[ht]
  \centering\small
  \caption{The annual loss rate per unit of exposure in the event-timed pool (years 2 to 30 pooled)
  against the approximation used for the premium: mean loss; the 99.9\% quantile, event-timed; the
  replenished quantile of Proposition~\ref{prop:replenish} plus the granularity adjustment; and the
  approximation's error relative to the event-timed value.}
  \label{tab:eventq}
  \begin{tabular}{rrrrr}
    \toprule
    $N$ & mean loss & 99.9\%, event-timed & replenished $+$ GA & error \\
    \midrule
    \inputrows{data/table_event_quantile.tex}
    \bottomrule
  \end{tabular}
\end{table}

% ---------------------------------------------------------------------------------------------
\section{Limitations}\label{sec:limits}

\begin{enumerate}
\item One Gaussian factor. The quantiles are those of the Gaussian one-factor model at the stated
  correlation, and they are not bounds: correlation does not order all upper-tail quantiles across
  dependence structures. Adding independent Gaussian factors with the same total latent covariance
  leaves the distribution unchanged, while a factor with a heavier tail moves the far quantiles by
  much more than the correlation sweep of Table~\ref{tab:rhogrid} does; with a Student-$t$ factor at the
  same linear correlation the 99.9\% quantile is \TailTTenRatio\ to \TailTFiveRatio\ times the Gaussian
  one (Section~\ref{sec:tail}). The premiums of Section~\ref{sec:pricing} are conditional on the
  Gaussian factor, and a pool whose borrowers share a tail risk needs a dependence assumption of its own.
\item PD and correlation are scenarios. The Basel retail correlation comes from bank books; the grid up
  to 0.30 brackets more concentrated pools.
\item Independent years in the reserve simulation. Credit cycles are serially correlated, which makes
  long runs of losses, and so lender losses under a reserve at $\EL/a$, more likely than shown.
\item Constant deposits in the reserve simulation. Losses and withdrawals are replenished by new
  deposits, so exposure and the reserve's inflow never shrink; Section~\ref{sec:event} tests the
  recursion's timing, replacement and granularity assumptions in a finite pool but keeps this one.
\item A constant hazard across a borrower's loans, with no seasoning and no strategic default on the
  last loan.
\item Deterministic LGD, with no time value on the 60 days between disbursement and write-off.
\item Equal exposures; concentration beyond the granularity adjustment is not modelled, which the
  pool's maximum loan size limits.
\item Simple yields, a fixed funding rate and exogenous utilisation. In stress, withdrawals would move
  utilisation.
\item Not modelled: fraud by the issuer of credit lines, smart-contract failure, a loss of the
  stablecoin's peg, and loans repaid within 24 hours, which pay no interest yet carry some risk.
\end{enumerate}

% ---------------------------------------------------------------------------------------------
\section{Conclusion}

Pricing uncollateralized credit in an on-chain pool is standard portfolio credit risk once three
corrections are made: convert per-loan default rates to annual ones, account for re-lending, and
correct the tail for the pool's small size. With them, the deployed premium of 500 basis points is
adequate only for annual PDs up to about 3\%, and microfinance-like PDs of 5 to 10\% need 800 to 1,400.
The reserve is where the protocol's rules matter. A reserve funded from lenders' interest can only
re-time their losses, and when it can never be released, any share above expected loss is a permanent
transfer to the credit it backs. Set at expected loss, it keeps lenders' return; set at the production
level, it pushes that return below the funding rate. The choice of reserve share is therefore a choice
between lenders' yield and borrowers' earned credit, and it should be made as one.

\paragraph{Disclosure of AI authorship.} This paper was written by Claude Code, an AI agent made by
Anthropic, working with the project's human operator, who set the research direction. The calibration
code was written by AI agents working with the project; the simulation of the locked reserve is this
paper's own. All of it is public and reproducible.

{\sloppy\paragraph{Data and code availability.} The calibration code and results, at commit b725a85,
are in \texttt{analysis/credit\_risk} of
\url{https://github.com/scottonchain/microcredit-contract}. The reserve simulation and the scripts
that extract every table and figure are in \url{https://github.com/scottonchain/microcredit-theory}.\par}

\bibliographystyle{apalike}
\bibliography{references}

\appendix
\section{Closed Forms}\label{app:closed}

With $c=\Phi^{-1}(\mathrm{PD})$ and $\Phi_2(\cdot,\cdot;\varrho)$ the bivariate standard normal
distribution function with correlation $\varrho$, the Vasicek default rate $D$ satisfies
\begin{align*}
  \mathrm{Var}\,D&=\Phi_2(c,c;\rho)-\mathrm{PD}^2,\\
  \mathrm{ES}_q&=\mathrm{LGD}\cdot\frac{\mathrm{PD}-\Phi_2\bigl(c,\Phi^{-1}(q);-\sqrt{\rho}\bigr)}{1-q},\\
  \E\bigl[(D-k)^+\bigr]&=\mathrm{PD}-\Phi_2\bigl(c,y_k;-\sqrt{\rho}\bigr)-k\bigl(1-\Phi(y_k)\bigr),\qquad
  y_k=\frac{\sqrt{1-\rho}\,\Phi^{-1}(k)-c}{\sqrt{\rho}}.
\end{align*}
The first-order granularity adjustment to the $q$-quantile for $N$ equal loans is
\[
  \mathrm{GA}=\mathrm{LGD}\cdot\frac{\mu(1-\mu)(y-gb)-(1-2\mu)\,\varphi(g)\,b}{2N\varphi(g)\,b},
\]
with $y=\Phi^{-1}(q)$, $b=\sqrt{\rho/(1-\rho)}$, $g=(c+\sqrt{\rho}\,y)/\sqrt{1-\rho}$ and $\mu=\Phi(g)$
\citep{martin2002,gordy2004}. Each is implemented in the contract repository's
\texttt{analysis/credit\_risk/vasicek.py} and tested against quadrature, numerical differentiation or
Monte Carlo there.

\end{document}
